% ECONOMICS_PROBLEM: {"id": "C78-227", "title": "Dynamic matching with congestion, waiting and entry", "jel_code": "C78", "source_id": "OP-0009"}
% !TeX program = xelatex
\documentclass[11pt,a4paper]{article}
\usepackage[margin=23mm]{geometry}
\usepackage{fontspec}
\setmainfont{Latin Modern Roman}
\usepackage{amsmath,amssymb,mathtools}
\usepackage{xcolor,enumitem,booktabs,longtable,array,xurl,hyperref}
\definecolor{muted}{HTML}{53636C}
\hypersetup{colorlinks=true,urlcolor=blue,linkcolor=blue}
\setlength{\parindent}{0pt}
\setlength{\parskip}{5pt}
\newcommand{\R}{\mathbb R}
\newcommand{\E}{\mathbb E}
\newcommand{\N}{\mathbb N}
\newcommand{\ind}{\mathbf 1}
\newcommand{\Var}{\operatorname{Var}}
\newcommand{\Cov}{\operatorname{Cov}}
\newcommand{\supp}{\operatorname{supp}}
\DeclareMathOperator*{\argmin}{arg\,min}
\DeclareMathOperator*{\argmax}{arg\,max}
\newcommand{\entrymeta}[3]{{\sffamily\small\color{muted}\textbf{\texttt{#1}}\quad|\quad #2\par #3\par}}
\newcommand{\Problemref}[1]{\texttt{#1}}
\newcommand{\JELref}[2]{\texttt{#2} (JEL \texttt{#1})}
\begin{document}
\section*{Dynamic matching with congestion, waiting and entry}
\textbf{Problem ID:} \texttt{C78-227}\quad \textbf{JEL:} \texttt{C78}\par
% BEGIN ECONOMICS STATEMENT
\label{problem:OP-0009}
\label{atlas:prob:M9}

\noindent\textbf{Original atlas ID:} M9. \textbf{Formulation:} editorial specification of a research family.

\subsection*{Definitions and assumptions}

Fix a finite horizon $t=0,\ldots,T$. A bounded finite set of potential agents with finite private type and action sets arrives according to a specified stochastic law; agents have private match preferences and waiting costs, and may leave according to a declared departure process. At each history $h_t$, a policy $m$ selects a feasible matching among available agents and specifies messages, priorities, and information disclosure. An agent's discounted utility is match value at the matching date, less discounted waiting costs, with a specified unmatched outside option. Let $\mathcal E(m)$ denote sequential-equilibrium assessments of this finite extensive-form market, incorporating reporting, acceptance, waiting, and entry decisions. When arrivals or departure hazards respond to payoffs, their behavioral laws are primitives or equilibrium objects that must be specified rather than treated as exogenous data.

\subsection*{Formal research question}

For a feasible policy class $\mathfrak M_F$, define $W(m,e)$ as expected discounted total match surplus minus waiting and administrative costs under $e\in\mathcal E(m)$. Investigate
\[
V^*=\sup_{m\in\mathfrak M_F}\inf_{e\in\mathcal E(m)}W(m,e).
\]
Characterize information and arrival-process restrictions under which waiting to increase market thickness improves welfare, and design policies with a proved bound $W(m,e)\ge V^*-\varepsilon$ for every selected admissible equilibrium. If dynamic stability is required, first define its deviation opportunities: a blocking pair must have a specified feasible time, information, and continuation alternative. Identify policy effects from randomized timing or disclosure while following later matches and exits, rather than conditioning only on agents remaining in the queue.

\subsection*{Interpretation and scope}

Static stability excludes an immediately profitable blocking pair in a fixed market. It does not determine optimal waiting with uncertain future partners, nor does it imply ex-ante efficiency. Congestion can mean queueing costs, limited capacity, or competition for partners; each has to appear in feasibility or utility. A fully observed market with exogenous arrivals may reduce to a stochastic control problem, whereas strategic reports and entry add incentive constraints and equilibrium selection. These are different benchmarks. Market-thickness experiments should measure changes in endogenous participation and attrition, because gains among surviving agents can coexist with losses for agents who depart unmatched. Infinite-horizon steady-state conclusions additionally require recurrence and integrability conditions.

\subsection*{Source audit and status}

Gale--Shapley (1962), Roth (1984), and Abdulkadiroğlu--Sönmez (2003) are verified. Original link [14] correctly identifies the school-choice paper, but these references mainly motivate static matching and institutional design; they do not directly source a general dynamic matching problem. Akbarpour--Li--Oveis Gharan (2020), added as [S4], studies stochastic arrivals/departures, thickness, and information, and already provides results in a concrete dynamic benchmark. The entry is an editorial extension to strategic congestion, timing, and entry. Its broad remaining agenda is not evidence that the cited static stability theorem or the 2020 waiting benchmark is unresolved.

\subsection*{References}

\begin{enumerate}

\item[S1] David Gale and Lloyd S. Shapley (1962). \emph{College Admissions and the Stability of Marriage}. American Mathematical Monthly 69(1), 9--15. \url{https://www.tandfonline.com/doi/abs/10.1080/00029890.1962.11989827}. \textit{Locator:} stable-matching construction. \textit{Establishes:} Existence and construction of stable matchings in the stated static model. \textit{Does not establish:} optimal waiting or entry in dynamic markets.

\item[S2] Alvin E. Roth (1984). \emph{The Evolution of the Labor Market for Medical Interns and Residents: A Case Study in Game Theory}. Journal of Political Economy 92(6), 991--1016. \url{https://stanford.edu/~alroth/evolut.html}. \textit{Locator:} sections 2--3. \textit{Establishes:} An institutional case study connecting clearinghouse performance and matching stability. \textit{Does not establish:} a general dynamic optimal mechanism.

\item[S3] Atila Abdulkadiroğlu and Tayfun Sönmez (2003). \emph{School Choice: A Mechanism Design Approach}. American Economic Review 93(3), 729--747. \url{https://www.aeaweb.org/articles?id=10.1257/000282803322157061}. \textit{Locator:} abstract; school-choice mechanisms. \textit{Establishes:} Static allocation design using priorities and strategic considerations. \textit{Does not establish:} dynamic waiting, congestion, and endogenous entry theory.

\item[S4] Mohammad Akbarpour, Shengwu Li, and Shayan Oveis Gharan (2020). \emph{Thickness and Information in Dynamic Matching Markets}. Journal of Political Economy 128(3), 783--815. \url{https://www.journals.uchicago.edu/doi/10.1086/704761}. \textit{Locator:} abstract; stochastic arrival/departure benchmark. \textit{Establishes:} Dynamic matching gains from waiting and information in a specified stochastic model. \textit{Does not establish:} optimal design for every strategic congestion market.

\end{enumerate}

\subsection*{Common conventions: Source mathematical conventions}

Unless an entry states otherwise, agent and firm sets are finite. State and action spaces are measurable, strategies and outcome maps are measurable, probability laws are specified, and expectations exist for the targets under discussion. Infinite-horizon discount factors lie in $(0,1)$ and discounted objectives are finite. Norms, tolerances and complexity input sizes must be specified for computational comparisons. Constants or primitives that appear locally are inputs, not empirical facts asserted by this catalogue.

\textbf{Identification.} A statistical model $m\in\mathcal M$ implies an observable law $P_m$ and target $\theta(m)$. At law $P$, its identified set is
\[
\Theta(P;\mathcal M)=\{\theta(m):m\in\mathcal M,\ P_m=P\}.
\]
Point identification holds when this set is a singleton, or equivalently when $P_{m_1}=P_{m_2}$ implies $\theta(m_1)=\theta(m_2)$ for all compatible models. Sharp bounds mean bounds attained, or approached when the boundary is not attained, by compatible models. Writing this set is a definition, not a solution: the substantive task is to establish informative restrictions and derive or estimate the set. An unrestricted-confounding model may give only trivial bounds.

\textbf{Inference.} Identification, estimability and valid uncertainty quantification are separate questions. Fix $\alpha\in(0,1)$. A confidence procedure $C_n$ over a declared law class $\mathcal P$ has asymptotic uniform coverage at level $1-\alpha$ when
\[
\liminf_{n\to\infty}\inf_{P\in\mathcal P}\Pr_P\{\theta(P)\in C_n\}\geq1-\alpha.
\]
For set-identified targets the coverage event must be defined separately, for example coverage of the full identified set. If an entry uses identification-indexed models instead of uniquely indexed laws, the same coverage convention is applied over those models. Pointwise limits do not establish uniform validity, and asymptotic validity does not guarantee finite-sample precision. Sample sizes, dependence structures and weak-identification sequences are part of each application.

\textbf{Counterfactuals.} $Y_i(a)$ is a potential outcome under a specified intervention, including its timing, implementation and versions. It is not $Y_i$ conditional on voluntarily choosing $a$. A full intervention vector permits interference; shorthand scalar treatment notation does not silently impose no interference. Assignment, selection, attrition, anticipation and measurement must be specified. A claim about an instrument's local effect is not automatically a population policy effect.

\textbf{Economic equilibrium.} A counterfactual model includes best responses, constraints, feasibility, market clearing, state transitions and expectations consistent with the stated information. When several equilibria exist, either a declared selector is an input or the target is set-valued. Comparisons across policy regimes must use the stated selection convention. Reduced-form relationships are not presumed invariant after an equilibrium-changing intervention.

\textbf{Optimization and welfare.} Optimization is over the stated feasible set. If attainment is not established, $\sup$ or $\inf$ and $\varepsilon$-optimal choices replace an asserted maximizer or minimizer. For example, with value $V=\sup_{a\in\mathcal A}W(a)<\infty$, an $\varepsilon$-optimal set is $\{a\in\mathcal A:W(a)\geq V-\varepsilon\}$ for $\varepsilon>0$. Benefits, costs, risk and health or environmental outcomes can be aggregated only using declared units and conversion weights. Interpersonal utility weights, the treatment of fiscal transfers and foreign profits, discounting, and the behavioral welfare criterion are explicit inputs. A comparative-static derivative additionally requires existence and appropriate differentiability of the selected counterfactual.

\textbf{Sources.} Local labels [S1], [S2], and so forth restart in every question. Locators identify the relevant abstract, model, section or result at the level actually checked; a bibliographic or abstract-level verification is not represented as inspection of an entire proof. Working papers are distinguished from later publications and changing versions. Source summaries are paraphrased. All URL checks and editorial formulations refer to the audit date stated above.
% END ECONOMICS STATEMENT
\end{document}
